\VerdierSemanticAnchor{front-title}
\VerdierPhysicalPageAnchor{0002}
\begin{titlepage}
  \centering
  {\Huge\bfseries 239\hfill ASTÉRISQUE\par}
  \vspace{1.5em}
  {\LARGE\bfseries 1996\par}
  \vfill
  {\Huge DES CATÉGORIES DÉRIVÉES\par}
  \vspace{0.4em}
  {\Huge DES CATÉGORIES ABÉLIENNES\par}
  \vspace{3em}
  {\Large\bfseries Jean-Louis VERDIER\par}
  \vfill
  {\Large\bfseries SOCIÉTÉ MATHÉMATIQUE DE FRANCE\par}
  \smallskip
  {\small\bfseries Publié avec le concours du CENTRE NATIONAL DE LA RECHERCHE SCIENTIFIQUE\par}
\end{titlepage}

\VerdierSemanticAnchor{front-ams-classification}
\VerdierPhysicalPageAnchor{0003}
\chapter*{A.M.S. Subjects classification}
\addcontentsline{toc}{chapter}{A.M.S. Subjects classification}

14F20, 18E30, 18E35, 18G05, 18G10, 18G15, 18G20,18G35,
18G40, 55U15, 55U25, 55U30.

\clearpage
\VerdierSemanticAnchor{front-portrait}
\VerdierPhysicalPageAnchor{0004}
\thispagestyle{empty}
\vspace*{0.62\textheight}
\begin{center}
Jean-Louis VERDIER, 1977
\end{center}

\clearpage
\VerdierPhysicalPageAnchor{0005}
\thispagestyle{empty}
\null
\clearpage

\VerdierSemanticAnchor{front-preface}
\VerdierPhysicalPageAnchor{0006}
\chapter*{Préface}
\addcontentsline{toc}{chapter}{Préface}

Jean-Louis Verdier, tragiquement disparu le 25 août 1989, a soutenu sa thèse
de doctorat d'État, intitulée \emph{Des catégories dérivées des catégories
abéliennes}, le 14 juin 1967, à Paris. Le présent texte est celui du manuscrit
dactylographié de cette thèse, dans sa dernière version. Les éditeurs se sont
bornés à corriger des erreurs de détail. Seule la partie intitulée
\emph{Introduction} avait été déposée à la faculté. Le corps du texte était
inachevé à la date de la soutenance, et l'est resté.

La genèse de ce travail est assez bien connue (\emph{cf.} [I]). On sait que
Grothendieck conçut l'idée des catégories dérivées au début des années soixante
pour fournir le cadre d'algèbre homologique nécessaire à la vaste
généralisation des théorèmes de dualité qu'il avait imaginée à la suite de son
exposé au congrès international d'Édimbourg de 1958. Il avait proposé à
Verdier, comme sujet de thèse, de construire le formalisme envisagé. En 1963,
Verdier faisait paraître un fascicule de résultats résumant l'essentiel de la
théorie, \emph{Catégories dérivées, quelques résultats (État 0)}, miméographié
par l'IHÉS (reproduit dans [SGA 4 1/2]). Si la notion de catégorie dérivée est
due à Grothendieck, c'est à Verdier que revient d'avoir introduit, en amont,
celle de catégorie triangulée. L'axiomatique présentée là, inspirée de la
théorie homotopique stable (\emph{cf.} [P]), devait se révéler d'une surprenante
fécondité.

Ce n'est que plusieurs années après que Verdier entreprend une rédaction
d'ensemble de son travail. Pourquoi ne l'a-t-il pas terminée? Bien des
hypothèses ont été avancées. Je pencherais pour la suivante. La partie qui
manque est la théorie des foncteurs dérivés. Quand Verdier en arrive à cet
endroit de sa rédaction, les catégories dérivées ont déjà été largement
utilisées en géométrie algébrique. Diverses variantes et généralisations se
développent : notion de dérivabilité ponctuelle à la Deligne, foncteurs dérivés
non additifs (algèbre homotopique à la Quillen), foncteurs dérivés filtrés. Dès
lors, le cadre fixé par le fascicule de résultats s'avérait insuffisant.
Trouver le bon point de vue aurait demandé un nouveau travail de fondements
que Verdier, intéressé à cette époque par d'autres problèmes, n'a eu, me
semble-t-il, ni le loisir ni sans doute l'envie de réaliser. À mesure que le
temps passait, des articles sur les catégories dérivées paraissaient et l'usage
s'en répandait, le formalisme envahissant d'autres domaines des mathématiques
(comme les équations aux dérivées partielles linéaires). L'achèvement de son
ouvrage devenait alors aux yeux de Verdier une tâche de moins en moins
nécessaire \ldots{}

\VerdierPhysicalPageAnchor{0007}
Tout incomplet qu'il est, ce texte de Verdier constitue néanmoins une
référence précieuse, voire irremplaçable (par exemple, la présentation,
originale et élégante, qui y est donnée des suites spectrales ne se trouve, à
ma connaissance, nulle part ailleurs). On ne peut que se réjouir qu'il voie
enfin le jour.

\begin{flushright}
Paris, le 27 avril 1996\\
Luc Illusie
\end{flushright}

\section*{Références}

\begin{description}
\item[{[I]}] L. Illusie, \emph{Catégories dérivées et dualité, travaux de
J.-L. Verdier}, L'Enseignement Mathématique, 36 (1990), 369-391. (Exposé donné
le 19 octobre 1989, lors de la cérémonie en hommage à Jean-Louis Verdier
organisée par l'Université Paris VII.)

\item[{[P]}] D. Puppe, \emph{On the formal structure of stable homotopy
theory}, Colloquium on Algebraic Topology, Aarhus Universitet (1962), 65-71.

\item[{[SGA 4 1/2]}] ``Cohomologie étale'', Séminaire de Géométrie Algébrique
du Bois-Marie, SGA 4 1/2, par P. Deligne, Lecture Notes in Mathematics 569,
Springer-Verlag (1977), 262-311.
\end{description}

